\documentclass[11pt]{book}
\usepackage{amsmath,amssymb}
\usepackage{tcolorbox}

\begin{document}

\chapter{Probability}

\section*{1.1 The Probabilistic Method}
\begin{tcolorbox}
\section*{概率方法}
\end{tcolorbox}

This chapter introduces the method.

Definition 1.1.1. A graph $G$ is \emph{sparse} if it has few edges.
% Definition 9.9.9 this comment line must not be detected as a candidate
Here is the body of the definition, which continues for a while.

Theorem 2.3.4 (Erd\H{o}s-Szekeres). Every sequence of $n^2+1$ distinct real numbers
contains a monotone subsequence of length $n+1$.

Proof. Fix a sequence. We argue by pigeonhole.
\begin{align}
a^2 + b^2 &= c^2 \\
d^2 &= e^2
\end{align}
This completes the proof.

A known result by Lemma 3.1 shows that the bound is sharp.
\[
  \int_0^1 f(x)\,dx
\]

Remark. This is a short note without a number.

\begin{theorem}
Theorem 9.9. A statement.
\end{theorem}
The body of Theorem 9.9 was left outside the environment by mistake.

\begin{verbatim}
Theorem 0.0 inside verbatim must be ignored.
\end{verbatim}

\section*{EXERCISES}
\begin{tcolorbox}
\section*{练习}
\end{tcolorbox}
\begin{enumerate}
\item[1.] Prove that $1+1=2$.
\item[2.] Show that every graph is sparse.
\end{enumerate}

\section*{Real exercises}

1. First problem text with a formula $a+b=c$.

2. Second problem text.

3. Third problem text.
Show that the third problem follows from the first two.

\end{document}
